\subsubsection{绝育后悔：发生率、原因与决策前的自我评估}

绝育手术最需要正视的议题是\textbf{后悔（regret）}——它不是小概率的意外，而是有规律可循的已知现象，应当在术前被充分讨论，而不是术后被独自消化。

\vspace{0.5em}
\noindent\textbf{发生率与规律}：

\begin{itemize}
	\item 各国研究数据差异较大，但方向一致：\textbf{绝育后悔率随年龄增长而明显下降}——35 岁以下、子女数少或尚无子女者做决定时的后悔风险显著高于 35 岁以上已育人群
	\item \textbf{后悔的常见触发事件}：离异或丧偶后与新伴侣组成新家庭、子女意外离世（参见"失独家庭"一节）、家庭经济或观念的重大变化、以及\textbf{在关系不稳定时期做出的"赌气式"或"讨好式"决定}
\end{itemize}

\vspace{0.5em}
\noindent\textbf{复通的现实}：输精管复通术的再通成功率约 50\%--80\%，输卵管复通约 30\%--60\%——但\textbf{再通成功不等于能自然怀孕}，实际妊娠率更低，且手术费用高、需要显微外科条件；复通失败或不适合复通者，辅助生殖（IVF）是另一条路，同样费用不菲且无保证。\textbf{正确的预期是：把绝育当作不可逆的决定，把复通当作"最后的补救"而非"后悔药"}。有再生育顾虑的男性，可考虑术前冻存精液作为备份（参见男性卷生育力保存相关内容）。

\vspace{0.5em}
\noindent\textbf{决策前的自我评估清单}：

\begin{itemize}
	\item 如果明天伴侣关系结束，\textbf{我还想不想要（与别人生）孩子}？
	\item 我做这个决定，是自己的意愿，还是为了\textbf{讨好伴侣、回应催婚催育压力、或解决婚姻矛盾}？（关系动荡期不做永久决定）
	\item 如果几年后孩子意外离世或家庭变故，我能接受这个选择不可撤销吗？
	\item 长效可逆避孕（宫内节育器、皮下埋植剂——避孕效果接近绝育、取出即恢复生育）\textbf{为什么不能满足我}？
	\item 我是否了解：输精管结扎在创伤、费用与并发症上\textbf{显著低于}女性输卵管手术——若双方商量"一人去结扎"，选择权与理由是否讨论过？
\end{itemize}

\begin{tcolorbox}[colback=pink!5!white,colframe=red!50!black,title=贴心小叮咛]
绝育适合"非常确定"的人，而不适合"暂时想省事"的人。凡是犹豫，就选可逆的方案——避孕是长期的事，你永远有下一次选择的机会，前提是这一次不把自己锁死。
\end{tcolorbox}
